\documentclass{lecture-kuznia}
\usepackage{textcomp}
\usepackage{comment}

\def\title{Turnieje: porządek ukryty w chaosie}
\def\data{2026}
\def\place{Poręba Wielka}
\def\group{Finaliści}
\def\author{Autor: Stefan Świerczewski}
\def\lecturer{Prowadzący: Stefan Świerczewski}
\def\event{WM 2026}

\ifshowsolutions
  \newenvironment{solutionbox}{%
    \par\smallskip\begingroup\color{black!85}%
    \noindent\textbf{Rozwiązanie. }%
  }{%
    \hfill$\blacksquare$\par\smallskip\endgroup
  }
\else
  \excludecomment{solutionbox}
\fi

\newcounter{zadanie}
\newenvironment{zadanie}[1][]{%
  \refstepcounter{zadanie}%
  \par\medskip\noindent
  \textbf{Zadanie \thezadanie.}%
  \ifstrempty{#1}{}{ \textit{(#1)}}%
  \quad
}{\par\medskip}

\newcommand{\outn}[1]{N^+(#1)}
\newcommand{\inn}[1]{N^-(#1)}
\newcommand{\outd}[1]{d^+(#1)}

\begin{document}

\begin{center}
  {\Huge\textbf{\title}}

  \medskip
  {\large Dziewięć problemów o szukaniu struktury w grafie skierowanym}
\end{center}

\fontsize{12}{16}\selectfont

\section{Tylko język}

\textbf{Turniejem} nazywamy graf otrzymany przez skierowanie każdej krawędzi
grafu pełnego. Zapis $x\to y$ oznacza, że krawędź pomiędzy $x$ i $y$ jest
skierowana od $x$ do $y$. Można myśleć, że zawodnik $x$ wygrał z zawodnikiem
$y$; remisów nie ma i każda para rozegrała dokładnie jeden mecz.

Dla wierzchołka $v$ oznaczamy
\[
  \outn v=\{x:v\to x\},\qquad
  \inn v=\{x:x\to v\},\qquad
  \outd v=|\outn v|.
\]
Liczbę $\outd v$ nazywamy \textbf{wynikiem} wierzchołka $v$.

Turniej jest \textbf{silnie spójny}, jeśli z każdego wierzchołka można dojść
do każdego innego, idąc zgodnie ze strzałkami. Jest \textbf{przechodni}, jeśli
jego wierzchołki można ustawić w kolejności $v_1,\ldots,v_n$ tak, że
$v_i\to v_j$ dla wszystkich $i<j$. Wierzchołek $v$ jest \textbf{królem}, jeśli
do każdego innego wierzchołka można dojść z $v$ w co najwyżej dwóch krokach.

To cała teoria potrzebna na wejściu. W zadaniach będą stale wracały cztery
ruchy:
\begin{itemize}
  \item wybór wierzchołka o ekstremalnym wyniku;
  \item maksymalizacja ścieżki, cyklu albo uporządkowania;
  \item zliczanie tego samego obiektu na dwa sposoby;
  \item losowanie turnieju i oszacowanie wartości oczekiwanej.
\end{itemize}

\section{Pierwsza runda}

\begin{zadanie}[ścieżka zamiast tabeli wyników]
Udowodnij, że wierzchołki każdego turnieju można ustawić w kolejności
\[
  v_1,v_2,\ldots,v_n
\]
tak, aby $v_i\to v_{i+1}$ dla każdego $i<n$.

Udowodnij również, że turniej jest przechodni wtedy i tylko wtedy, gdy nie
zawiera skierowanego trójkąta.
\end{zadanie}

\begin{solutionbox}
Pierwszą część dowodzimy indukcyjnie. Usuńmy wierzchołek $x$ i ustawmy
pozostałe w ścieżkę
\[
  v_1\to v_2\to\ldots\to v_{n-1}.
\]
Jeżeli $x\to v_1$, wstawiamy $x$ na początku. Jeżeli $v_{n-1}\to x$,
wstawiamy go na końcu. W przeciwnym razie, idąc wzdłuż ścieżki, w pewnym
miejscu po raz pierwszy zmienia się relacja z $x$. Znajdziemy więc indeks $i$
taki, że
\[
  v_i\to x\to v_{i+1},
\]
i wstawimy $x$ pomiędzy $v_i$ i $v_{i+1}$.

Turniej przechodni nie zawiera skierowanego cyklu, a więc w szczególności
skierowanego trójkąta. W drugą stronę weźmy otrzymaną wyżej ścieżkę. Gdyby
dla pewnych $i<j$ zachodziło $v_j\to v_i$, wybierzmy taką parę z najmniejszą
różnicą $j-i$. Mamy $j-i>1$. Jeżeli $v_{i+1}\to v_j$, to
\[
  v_j\to v_i\to v_{i+1}\to v_j
\]
jest skierowanym trójkątem. W przeciwnym razie $v_j\to v_{i+1}$, wbrew
minimalności różnicy $j-i$. Zatem każda krawędź biegnie do przodu i turniej
jest przechodni.
\end{solutionbox}

\begin{zadanie}[król]
Udowodnij, że każdy turniej ma króla.
\end{zadanie}

\begin{solutionbox}
Wybierzmy wierzchołek $v$ o największym wyniku. Rozważmy dowolny wierzchołek
$u$, dla którego $u\to v$. Gdyby nie istniał wierzchołek $w$ spełniający
$v\to w\to u$, to $u$ wygrywałby ze wszystkimi zawodnikami pokonanymi przez
$v$, a dodatkowo także z samym $v$. Wtedy
\[
  \outd u\geq \outd v+1,
\]
co przeczy wyborowi $v$. Każdy zawodnik jest więc osiągalny z $v$ w jednym
albo dwóch krokach.
\end{solutionbox}

\begin{zadanie}[cykl przez wszystkich]
Udowodnij, że każdy silnie spójny turniej na co najmniej trzech wierzchołkach
ma skierowany cykl przechodzący przez wszystkie wierzchołki dokładnie raz.
\end{zadanie}

\begin{solutionbox}
Weźmy najdłuższy skierowany cykl $C$. Ma on co najmniej trzy wierzchołki:
najkrótszy skierowany cykl w turnieju musi być trójkątem, bo każda cięciwa
dłuższego cyklu skracałaby go z jednej albo z drugiej strony.

Przypuśćmy, że $x\notin C$. Jeżeli $x$ zarówno wygrywa, jak i przegrywa z
wierzchołkami $C$, to na cyklu występują dwa kolejne wierzchołki $c_i,c_{i+1}$
takie, że
\[
  c_i\to x\to c_{i+1}.
\]
Można wtedy wstawić $x$ do $C$, sprzecznie z maksymalnością. Każdy wierzchołek
spoza $C$ należy zatem do jednego z dwóch zbiorów
\[
  A=\{x:C\to x\},\qquad B=\{x:x\to C\},
\]
gdzie $C\to x$ oznacza, że każdy wierzchołek cyklu wygrywa z $x$.

Silna spójność wymusza $A\neq\varnothing$ i $B\neq\varnothing$. Ponadto nie
wszystkie krawędzie pomiędzy $A$ i $B$ mogą biec z $B$ do $A$, bo wtedy z
$A$ nie dałoby się dojść do $B$ ani do $C$. Istnieją więc $a\in A$ i $b\in B$
takie, że $a\to b$. Dowolną krawędź $c_i\to c_{i+1}$ cyklu można zastąpić
ścieżką
\[
  c_i\to a\to b\to c_{i+1},
\]
otrzymując dłuższy cykl. Sprzeczność dowodzi, że $C$ zawiera wszystkie
wierzchołki.
\end{solutionbox}

\section{Zliczanie i przypadek losowy}

\begin{zadanie}[ile może być paradoksów?]
Trójkę wierzchołków nazwiemy \textit{cykliczną}, jeśli tworzy ona skierowany
trójkąt. Wyznacz największą możliwą liczbę cyklicznych trójek w turnieju na
$n$ wierzchołkach. Scharakteryzuj przypadek równości za pomocą wyników
wierzchołków.
\end{zadanie}

\begin{solutionbox}
Niech wyniki wierzchołków będą równe $d_1,\ldots,d_n$. Każda trójka jest
albo cykliczna, albo przechodnia. W trójce przechodniej istnieje dokładnie
jeden wierzchołek wygrywający z pozostałymi dwoma. Wobec tego liczba trójek
przechodnich wynosi
\[
  \sum_{i=1}^n\binom{d_i}{2}.
\]
Jeśli $T$ oznacza liczbę trójek cyklicznych, to
\[
  T=\binom n3-\sum_{i=1}^n\binom{d_i}{2}.             \tag{1}
\]
Ponadto $\sum_i d_i=\binom n2$. Dla ustalonej sumy wyrażenie
$\sum_i\binom{d_i}{2}$ jest najmniejsze, gdy liczby $d_i$ różnią się o co
najwyżej $1$: jeśli $a\geq b+2$, to zastąpienie pary $(a,b)$ przez
$(a-1,b+1)$ zmniejsza tę sumę.

Jeżeli $n$ jest nieparzyste, dostajemy $d_i=(n-1)/2$ dla każdego $i$ i
\[
  T_{\max}=\frac{n(n^2-1)}{24}.
\]
Jeżeli $n$ jest parzyste, połowa wyników jest równa $n/2$, a połowa
$n/2-1$, skąd
\[
  T_{\max}=\frac{n(n^2-4)}{24}.
\]

Równość rzeczywiście jest osiągalna. Dla $n=2m+1$ ustawmy wierzchołki
cyklicznie modulo $n$ i skierujmy $i\to j$, gdy
$j-i\pmod n\in\{1,\ldots,m\}$. Każdy wynik jest równy $m$. Dla parzystego
$n$ usuwamy jeden wierzchołek z takiego regularnego turnieju na $n+1$
wierzchołkach. Wzór (1) pokazuje zarazem, że opisane rozkłady wyników są
konieczne dla równości.
\end{solutionbox}

\begin{zadanie}[duży uporządkowany fragment]
Niech $k\geq 1$.
\begin{enumerate}[label=\alph*)]
  \item Udowodnij, że każdy turniej na co najmniej $2^{k-1}$ wierzchołkach
  zawiera przechodni podturniej na $k$ wierzchołkach.
  \item Udowodnij, że dla każdych $k\leq N<2^{(k-1)/2}$ istnieje turniej na
  $N$ wierzchołkach, który nie zawiera przechodniego podturnieju na $k$
  wierzchołkach.
\end{enumerate}
\end{zadanie}

\begin{solutionbox}
Pierwszą część dowodzimy indukcyjnie po $k$. Wybierzmy dowolny wierzchołek
$v$. Co najmniej jeden ze zbiorów $\outn v,\inn v$ ma co najmniej
$2^{k-2}$ elementów. Z założenia indukcyjnego zawiera on przechodni
podturniej na $k-1$ wierzchołkach. Dołączając $v$ na początku albo na końcu
jego przechodniego porządku, otrzymujemy szukane $k$ wierzchołków.

W drugiej części skierujmy każdą krawędź niezależnie, z prawdopodobieństwem
$1/2$ w każdą stronę. Ustalony zbiór $k$ wierzchołków jest przechodni z
prawdopodobieństwem
\[
  \frac{k!}{2^{\binom k2}},
\]
bo każdemu z $k!$ liniowych porządków odpowiada dokładnie jedno przechodnie
skierowanie. Oczekiwana liczba przechodnich $k$-elementowych podturniejów
jest więc mniejsza lub równa
\[
  \binom Nk\frac{k!}{2^{\binom k2}}
  \leq \frac{N^k}{2^{k(k-1)/2}}<1.
\]
Istnieje zatem wybór kierunków, dla którego ta liczba jest równa zeru.
\end{solutionbox}

\begin{zadanie}[mała rada nadzorcza]
Zbiór $S$ wierzchołków turnieju nazwiemy dominującym, jeśli dla każdego
$x\notin S$ istnieje $s\in S$ takie, że $s\to x$. Udowodnij, że każdy
turniej na $n$ wierzchołkach ma zbiór dominujący o rozmiarze nie większym
niż
\[
  \left\lceil\log_2(n+1)\right\rceil.
\]
\end{zadanie}

\begin{solutionbox}
Wśród jeszcze nierozpatrzonych wierzchołków wybieramy wierzchołek $v$ o
największym wyniku w indukowanym podturnieju. Jeżeli pozostało $m$
wierzchołków, to $v$ wygrywa z co najmniej $(m-1)/2$ z nich. Dodajemy $v$ do
$S$ i usuwamy $v$ wraz ze wszystkimi pokonanymi przez niego wierzchołkami.

Jeżeli po $t$ krokach pozostało $m_t$ wierzchołków, to
\[
  m_{t+1}+1\leq\frac{m_t+1}{2},
\]
a więc $m_t+1\leq(n+1)/2^t$. Po
$t=\lceil\log_2(n+1)\rceil$ krokach nic już nie pozostaje. Każdy usunięty
wierzchołek spoza $S$ przegrał z wierzchołkiem dodanym do $S$, więc $S$ jest
dominujący.
\end{solutionbox}

\section{Finałowe rundy}

\begin{zadanie}[najlepsza kolejność wie więcej, niż powinna]
Udowodnij, że wierzchołki każdego turnieju można ustawić w kolejności
$v_1,\ldots,v_n$ o następujących własnościach: dla wszystkich $i<j$
\[
\begin{aligned}
  \big|\{k:i<k\leq j,\ v_i\to v_k\}\big|&\geq\frac{j-i}{2},\\
  \big|\{k:i\leq k<j,\ v_k\to v_j\}\big|&\geq\frac{j-i}{2}.
\end{aligned}                                                  \tag{2}
\]
Wywnioskuj z tego ponownie, że każdy turniej ma ścieżkę przechodzącą przez
wszystkie wierzchołki oraz króla.
\end{zadanie}

\begin{solutionbox}
Spośród wszystkich ustawień wybierzmy takie, w którym liczba krawędzi
biegnących do przodu, czyli krawędzi $v_p\to v_q$ z $p<q$, jest największa.

Gdyby pierwsza nierówność (2) nie zachodziła dla pewnych $i<j$, to wśród
$v_{i+1},\ldots,v_j$ więcej krawędzi biegłoby do $v_i$ niż od $v_i$.
Przeniesienie $v_i$ tuż za $v_j$ zwiększyłoby liczbę krawędzi biegnących do
przodu. Analogicznie, jeżeli nie zachodziłaby druga nierówność, przeniesienie
$v_j$ tuż przed $v_i$ poprawiłoby ustawienie. Obie sytuacje przeczą
maksymalności.

Dla $j=i+1$ pierwsza nierówność wymusza $v_i\to v_{i+1}$, więc otrzymaliśmy
ścieżkę Hamiltona.

Pokażemy, że $v_1$ jest królem. Jeśli $v_1\to v_j$, sprawa jest jasna.
Załóżmy więc, że $v_j\to v_1$. Z pierwszej nierówności dla przedziału
$v_1,\ldots,v_j$ wynika, że $v_1$ wygrywa z co najmniej
$\lceil(j-1)/2\rceil$ wierzchołkami spośród $v_2,\ldots,v_{j-1}$. Z drugiej
wynika, że co najmniej $\lceil(j-1)/2\rceil$ wierzchołków z tego samego
zbioru wygrywa z $v_j$. Zbiory te muszą się przeciąć, bo razem mają więcej
elementów niż zbiór $\{v_2,\ldots,v_{j-1}\}$. Istnieje zatem $k$ takie, że
\[
  v_1\to v_k\to v_j.
\]
\end{solutionbox}

\begin{zadanie}[gra w odwracanie trójkątów]
Na ustalonym, ponumerowanym zbiorze wierzchołków rozważamy dwa turnieje
$T$ i $T'$. Dozwolony ruch polega na wybraniu skierowanego trójkąta i
odwróceniu wszystkich jego krawędzi.

Udowodnij, że z $T$ można otrzymać $T'$ za pomocą takich ruchów wtedy i
tylko wtedy, gdy każdy wierzchołek ma ten sam wynik w $T$ i w $T'$.
\end{zadanie}

\begin{solutionbox}
Każdy wierzchołek skierowanego trójkąta ma w nim jedno zwycięstwo i jedną
porażkę. Odwrócenie trójkąta nie zmienia więc żadnego wyniku. Pozostaje
udowodnić drugą implikację.

Najpierw zauważmy lemat: \textit{krawędzie dowolnego skierowanego cyklu
można odwrócić za pomocą ruchów na trójkątach, nie zmieniając ostatecznie
żadnej innej krawędzi}. Dowód prowadzimy indukcyjnie po długości cyklu
\[
  x_1\to x_2\to\ldots\to x_r\to x_1.
\]
Dla $r=3$ nie ma czego dowodzić. Jeżeli $x_3\to x_1$, najpierw odwracamy
trójkąt $x_1x_2x_3$, a następnie krótszy cykl
\[
  x_1\to x_3\to x_4\to\ldots\to x_r\to x_1.
\]
Jeżeli $x_1\to x_3$, wykonujemy te same dwie operacje w odwrotnej kolejności.
Cięciwa $x_1x_3$ zostaje odwrócona dwa razy, a każda krawędź pierwotnego
cyklu dokładnie raz.

Zbudujmy teraz graf skierowany $D$: umieszczamy w nim krawędź $x\to y$, gdy
w $T$ mamy $x\to y$, a w $T'$ ta sama krawędź jest skierowana przeciwnie.
Równość wyników w $T$ i $T'$ oznacza, że w $D$ stopień wychodzący każdego
wierzchołka jest równy jego stopniowi wchodzącemu. Krawędzie $D$ można więc
rozłożyć na skierowane cykle. Odwracamy kolejno te cykle, korzystając z
lematu. Każda taka operacja uzgadnia z $T'$ wszystkie krawędzie danego cyklu
i nie narusza pozostałych. Po skończeniu otrzymujemy $T'$.
\end{solutionbox}

\begin{zadanie}[boss: parzystość, której nie widać]
Udowodnij, że liczba skierowanych ścieżek przechodzących przez wszystkie
wierzchołki turnieju dokładnie raz jest zawsze nieparzysta.
\end{zadanie}

\begin{solutionbox}
Oznaczmy zbiór krawędzi turnieju przez $E$. Dla permutacji
$\pi=(v_1,\ldots,v_n)$ niech
\[
  A_\pi=\{v_i\to v_{i+1}\in E:1\leq i<n\}.
\]
Dla $X\subseteq E$ oznaczmy przez $f(X)$ liczbę permutacji $\pi$ takich, że
$A_\pi=X$, a przez $g(X)$ liczbę permutacji takich, że $X\subseteq A_\pi$.
Mamy
\[
  g(X)=\sum_{Y\supseteq X}f(Y).
\]
Z zasady włączeń i wyłączeń otrzymujemy odwrócenie tej zależności:
\[
  f(X)=\sum_{Y\supseteq X}(-1)^{|Y|-|X|}g(Y).           \tag{3}
\]

Wartość $g(Y)$ jest niezerowa tylko wtedy, gdy krawędzie $Y$ tworzą rozłączne
skierowane ścieżki. Jeżeli wraz z izolowanymi wierzchołkami jest to $r$
ścieżek, każdą z nich traktujemy jak jeden blok i dostajemy $g(Y)=r!$.
Liczba ta jest nieparzysta dokładnie dla $r=1$, czyli wtedy, gdy $Y$ samo
jest ścieżką przechodzącą przez wszystkie wierzchołki. Z (3), po przejściu
modulo $2$, wynika zatem
\[
  f(X)\equiv h(X)\pmod 2,                               \tag{4}
\]
gdzie $h(X)$ oznacza liczbę ścieżek Hamiltona zawierających wszystkie
krawędzie z $X$.

Odwróćmy teraz jedną krawędź $a$ turnieju, otrzymując turniej $T'$. Tracimy
$h(\{a\})$ ścieżek Hamiltona. Ścieżki, które zyskujemy, po zapisaniu ich
w odwrotnej kolejności odpowiadają dokładnie tym permutacjom $\pi$ w $T$,
dla których $A_\pi=\{a\}$. Zyskujemy więc $f(\{a\})$ ścieżek. Na mocy (4)
obie liczby mają tę samą parzystość, zatem odwrócenie jednej krawędzi nie
zmienia parzystości liczby ścieżek Hamiltona.

Każdy turniej można otrzymać z turnieju przechodniego przez odwracanie
pojedynczych krawędzi. Turniej przechodni ma dokładnie jedną ścieżkę
Hamiltona. Każdy turniej ma ich więc nieparzyście wiele.
\end{solutionbox}

\ifshowsolutions
\section{Notatka dla prowadzącego}

Naturalny podział na dwa spotkania to zadania 1--6 oraz 7--9. Jeśli czasu jest
mniej, warto zrobić wspólnie zadania 2, 4 i 7, a następnie oddać grupie
zadanie 8. Zadanie 9 dobrze działa jako finał wykładu: potrzebuje jednego
nowego narzędzia, ale używa go do bardzo konkretnej obserwacji o parzystości.

Twierdzenia z zadań 1 i 9 są klasycznie przypisywane Rédeiemu, a twierdzenie
z zadania 3 --- Camionowi. Zadanie 7 wykorzystuje porządek medianowy. Dobór
trudności został porównany z zadaniami zawodów III stopnia Olimpiady
Matematycznej; materiał nie zakłada jednak znajomości żadnego twierdzenia o
turniejach.
\fi

\end{document}
